\section{Finite normalized simplicial input}
\label{sec:normalized-input}

Version 1.2.1 adds a finite input layer whose indices enumerate only
nondegenerate simplices.  The function \code{simplexCount n} fixes their
order in dimension $n$.  For a nondegenerate $(n+1)$-simplex $\sigma$, the
face table has type
\[
  \code{face}\ n :
  \operatorname{Fin}(n+2)\to
  \operatorname{Fin}(\code{simplexCount}(n+1))\to
  \operatorname{Option}
    (\operatorname{Fin}(\code{simplexCount}(n))).
\]
A result \code{some $\tau$} names a nondegenerate face.  A result
\code{none} says that the raw face is degenerate and therefore vanishes in
normalized chains.

The boundary matrix is derived entrywise:
\[
  (D_{n+1})_{\tau,\sigma}
  =
  \sum_{i=0}^{n+1}
  \begin{cases}
    (-1)^i,&\code{face}\ n\ i\ \sigma=\code{some}\ \tau,\\
    0,&\text{otherwise}.
  \end{cases}
\]
This definition handles two details without auxiliary postprocessing.
Repeated faces remain distinct summands and can cancel by sign, while
degenerate faces contribute zero immediately.  The fields of
\code{FiniteNormalizedSimplicialData} include finite support and a proof
that these derived matrices satisfy $D_{n+1}D_{n+2}=0$.  Thus a supplied
face table cannot drift from a separately stored boundary matrix.

The conversion \code{toIntChainComplex} is definitionally rank- and
boundary-preserving, so all kernel-coordinate and homology theorems apply
without a second representation bridge.  A dimension-two constructor asks
for vertex, edge, and triangle face tables and only the nonempty equation
$D_1D_2=0$.

Three executable fixtures isolate the conventions.  A minimal circle has one
vertex and one nondegenerate edge; its equal endpoint faces occur with
opposite signs.  The normalized model $\Delta^2/\partial\Delta^2$ has one
vertex, no nondegenerate edges, and one nondegenerate two-simplex whose three
faces are all \code{none}.  The ordinary filled triangle has the exact
oriented matrices
\[
D_1=
\begin{pmatrix}
-1&-1&0\\
 1& 0&-1\\
 0& 1&1
\end{pmatrix},
\qquad
D_2=
\begin{pmatrix}1\\-1\\1\end{pmatrix}.
\]
Kernel reduction verifies their entries and $D_1D_2=0$.
